\chapter{Процедура оценки и критерии}
\label{ch:evaluation}
\index{оценка}
\index{критерии оценки}

\section{Общие принципы}

Используются \textbf{25 критериев} (C01--C25), каждый не более \textbf{3} баллов (уровни \textbf{0 / 1 / 2 / 3}; подробно --- \path{docs/criteria\_rubric.md}); номинальная сумма \textbf{raw} --- \textbf{75}. В первую очередь оцениваются \textbf{изменения конкурсанта}: код в \texttt{src\_solution/} (\texttt{abu/tcb} и \texttt{abu/other}), SBOM в \texttt{src\_solution/sbom/}, сопровождаемые тесты и отчётные артефакты. Критерии, связанные с оценкой \textbf{свойств учебных подсистем ЦР и Регулятора} как объектов заслуги участника, \textbf{не применяются}; вместо них введены сквозной автотест сценария ЦР--АБУ и отдельные критерии по решению в \texttt{src\_solution/}.

Итог оценки в автоматической части выводится как сумма \textbf{raw} по критериям до \textbf{75}. Автоматический подсчёт части баллов выполняется скриптом \path{scripts/evaluate_contest_score.py} (команда \texttt{make evaluate-score} в репозитории).

\section{Краткий перечень критериев}

\begin{longtable}{@{}p{0.12\textwidth}p{0.82\textwidth}@{}}
\toprule
\textbf{ID} & \textbf{Содержание} \\
\midrule
\endfirsthead
\toprule
\textbf{ID} & \textbf{Содержание} \\
\midrule
\endhead
C01--C04 & Все тесты репозитория (в т.\,ч. тесты решения), тесты безопасности, маркер \texttt{security}, журнал \\
C05--C07 & Примеры SGA/SBOM, успешная сертификация решения (пакет + ответ Регулятора) \\
C08 & Сквозной функциональный автотест \textbf{ЦР--АБУ} (основной сценарий миссии) \\
C09 & Оформление кода в \texttt{src\_solution/} (flake8, PEP8) \\
C10--C12 & Решение: журнал, зависимости, тесты на \texttt{src\_solution/} \\
C13--C15 & Раздел тестов безопасности в отчёте решения; \texttt{numpy} в SBOM решения (TCB vs OTHER); тесты \texttt{event\_log} + \texttt{src\_solution} \\
C16 & Покрытие ДВБ решения (\texttt{src\_solution/abu/tcb}) тестами (\texttt{pytest --cov=src\_solution.abu.tcb}) \\
C17 & Отчёт конкурсанта в каталоге \texttt{src\_solution/docs/} \\
C18--C19 & Решение: \textbf{security\_monitor}, \textbf{policies}; изоляция доменов и контроль запросов/ответов \\
C20--C22 & Экспертные критерии жюри (стоимость сертификации, архитектура, воспроизводимость) \\
C23--C25 & Размер доменов ДВБ, количество интерфейсов, проверяемость монитора безопасности \\
\bottomrule
\end{longtable}

\section{Подсказка для проектирования (C18--C19)}

Полное описание шаблона монитора безопасности, изоляции доменов, политик и учебного примера «Светофор» приведено в разделе~\ref{sec:monitor-traffic}. Разбор TARA в контексте задания --- в разделе~\ref{sec:tara-contest}. Локальный путь к ноутбуку:\\
\path{references/traffic_light_demo/cyberimmunity-traffic-lights-example.ipynb}

\section{Отчёт конкурсанта и шаблоны}

Текстовый отчёт о решении оформляется в каталоге \path{src_solution/docs/} (критерий \textbf{C17}): обязательные разделы --- архитектура, политики безопасности, результаты сквозных и security-тестов, сертификация, архитектурные диаграммы или ссылки на них (подробная расшифровка --- приложение~\ref{app:criteria-rubric}). В репозитории приведён заполняемый шаблон отчёта. Для прозрачной сдачи дополнительно используется табличный шаблон баллов \path{docs/templates/evaluation_report.md} (см. приложение~\ref{app:template-appeal}).

\section{Регламент: разрешение равенства баллов и прозрачность}

Правила разрешения равенства по сумме баллов, требования к \textbf{прозрачности} сдачи и апелляции изложены в \textbf{полном регламенте} (приложение~\ref{app:regulations}, разделы «Разрешение равенства баллов» и «Прозрачность»). Шаблоны отчёта (в т.\,ч. таблица критериев) и порядок апелляции --- приложение~\ref{app:template-appeal}. Справочно по содержанию основного текста: сертификация --- глава~\ref{ch:certification}; инфраструктура и тесты --- глава~\ref{ch:infra}.

\section{Расшифровка уровней (0--3) по критериям}

Подробные формулировки уровней баллов приведены в приложении~\ref{app:criteria-rubric} и в файле \path{docs/criteria\_rubric.md} репозитория.
